%% festival_shocks.tex
%% Festivals as Demand Shocks: a decade of weather-adjusted estimates for
%% Delhi's peak electricity demand, and their value for day-ahead forecasting.
%%
%% Single-file manuscript. Compile with:
%%     pdflatex festival_shocks
%%     bibtex   festival_shocks
%%     pdflatex festival_shocks
%%     pdflatex festival_shocks
%%
%% All numbers in this draft are the feasibility estimates recorded in
%% docs/paper_festival_effects.md (run 2026-10-29). Figures referenced below
%% must be regenerated from scripts/paper/*.py before submission; until then the
%% \placeholderfig macro draws a labelled box so the document still compiles.
%% See docs/paper/README.md for the pre-submission checklist.

\documentclass[11pt]{article}

\usepackage[utf8]{inputenc}
\usepackage[T1]{fontenc}
\usepackage[margin=1in]{geometry}
\usepackage{amsmath,amssymb}
\usepackage{booktabs}
\usepackage{array}
\usepackage{graphicx}
\usepackage{siunitx}
\usepackage{xcolor}
\usepackage[numbers,sort&compress]{natbib}
\usepackage[hidelinks]{hyperref}

\graphicspath{{figures/}}

% siunitx: treat percentages cleanly, keep tabular numbers aligned
\sisetup{detect-weight=true, detect-family=true}

% --- placeholder-figure macro --------------------------------------------
% If the real figure file exists it is included; otherwise a labelled box is
% drawn so the manuscript compiles before the plots are regenerated.
\newcommand{\placeholderfig}[2]{%
  % #1 = file basename (no extension), resolved under figures/; #2 = height.
  % \IfFileExists ignores \graphicspath, so probe the figures/ path explicitly;
  % \includegraphics still resolves #1 via \graphicspath once the file exists.
  \IfFileExists{figures/#1.pdf}{\includegraphics[width=\linewidth]{#1}}{%
    \IfFileExists{figures/#1.png}{\includegraphics[width=\linewidth]{#1}}{%
      \setlength{\fboxsep}{0pt}%
      \fbox{\parbox[c][#2][c]{\dimexpr\linewidth-2\fboxrule\relax}{%
        \centering\ttfamily\small
        [\,figure placeholder: \detokenize{#1}\,]\\[2pt]
        \normalfont\itshape\footnotesize
        TODO: regenerate from scripts/paper/ (see README)%
      }}%
    }%
  }%
}

\newcommand{\pct}{\%}

\title{\textbf{Festivals as Demand Shocks:\\[2pt]
A Decade of Weather-Adjusted Estimates for Delhi's Peak\\
Electricity Demand, and Why Binary Holiday Flags\\
Under-serve Gradient-Boosted Forecasters}}

\author{Rupesh Bharambe\\
  \small Gridalytics\\
  \small \texttt{rupeshbharambe2004@gmail.com}}

\date{\today}

\begin{document}
\maketitle

\begin{abstract}
\noindent
Festivals are large, recurring, culturally-timed shocks to electricity demand,
yet published Indian work describes them only qualitatively. We give, to our
knowledge, the first multi-year (2015--2024) per-festival estimates of holiday
effects on Delhi's daily \emph{peak} demand. Each estimate is a weather-adjusted
counterfactual from an event-study regression with Newey--West HAC(7) standard
errors, validated by placebo-randomisation inference. The effects are large and
stable year to year: Holi suppresses the daily peak by about $31\pct$, Diwali by
about $17\pct$, and Republic Day by about $14\pct$, while a few festivals
(Christmas, Maha Shivaratri) raise it. The Diwali dip is a three-day structure
(Diwali, Govardhan Puja, Bhai Duj) and Holi carries a day-after effect, neither
of which a single binary holiday flag can represent. In a rolling-origin
day-ahead backtest (test years 2019--2024, LightGBM), encoding the festival as a
signed effect-size prior cuts festival-day MAPE from $6.6\pct$ to $5.5\pct$ at
horizon~1 and from $9.5\pct$ to $6.9\pct$ at horizon~7, whereas a binary flag or
a categorical festival identity do not: with roughly one occurrence per year, a
gradient-boosted tree cannot isolate a rare event from a dummy. We frame this as
a rare-event encoding result, not a new estimator. Finally, we present a
data-integrity case study: three silent pipeline defects each biased earlier
conclusions toward ``festivals don't matter,'' and a fourth---a target leak in
differenced/rolling features---flattered cross-validated accuracy; we give cheap
detection tests and report leak-free results after the fixes. The quantification
is robust and operationally relevant for load-dispatch scheduling and power
purchase; the forecasting lever is real but narrow, and we are deliberate about
where it is worth carrying.
\end{abstract}

\vspace{0.5em}
\noindent\textbf{Keywords:} electricity demand; peak load; festivals; holiday
effects; counterfactual estimation; event study; load forecasting; gradient
boosting; data integrity; India.

% =========================================================================
\section{Introduction}
\label{sec:intro}

Electricity demand in India swings sharply around festivals. Offices, markets and
industry close; domestic routines shift; and the diurnal load profile is
redrawn. These are among the largest short-horizon, non-weather deviations a load
forecaster faces, and they recur on culturally-timed dates---many of them lunar,
so they move through the Gregorian calendar from year to year. For a State Load
Dispatch Centre (SLDC) scheduling generation a day ahead, and for a distribution
company (DISCOM) sizing its day-ahead power purchase, mis-sizing a festival is a
concrete operational cost.

Despite this, the published Indian literature characterises festival effects only
\emph{qualitatively}. \citet{gaur2016} note an ``abrupt fall'' around festivals
but reuse weekend coefficients rather than estimating festival-specific
magnitudes; the Grid-India \emph{Electricity Demand Pattern Analysis}
\citep{gridindia2023} discusses festival behaviour narratively without per-festival
Delhi magnitudes; and \citet{hunt2025}, the nearest academic neighbour, model
Indian states but report no per-festival sizes. One study of the Northern grid
\citep{banerjee2010} found \emph{no} holiday effect at all.

This paper makes three empirical contributions and one methodological cautionary
tale, stated here in their narrowest honest form.

\begin{enumerate}
\item \textbf{Quantification.} We estimate, for each major festival, a
  weather-adjusted counterfactual effect on Delhi's daily peak demand over
  2015--2024, with uncertainty intervals (Table~\ref{tab:effects}). To our
  knowledge these are the first multi-year per-festival magnitudes for Delhi. The
  effects are large and year-to-year stable.
\item \textbf{Multi-day structure.} The Diwali effect spans about three days
  (Diwali, Govardhan Puja, Bhai Duj), and Holi has a pronounced day-after effect.
  A single-day binary flag cannot represent either pattern.
\item \textbf{Forecasting value of an effect-size prior.} In a rolling-origin
  day-ahead backtest, a causal signed-magnitude feature reduces festival-day
  error for a gradient-boosted model, while a binary flag and even a categorical
  festival identity do not. We frame this as a rare-event encoding result, not a
  new method.
\item \textbf{Data-integrity case study.} We document how several silent pipeline
  defects earlier manufactured a near-null result, and give the cheap tests that
  catch them.
\end{enumerate}

We deliberately \emph{do not} claim this to be the first such study for India at
large (state-level Grid-India tables are unchecked), nor that the mechanism is
new, nor any sub-daily accuracy headline: the quantification is daily peak only.

% =========================================================================
\section{Related work}
\label{sec:related}

\paragraph{Indian demand studies.} \citet{gaur2016} and the Grid-India demand
pattern reports \citep{gridindia2023} describe festival and weekend behaviour of
Indian load qualitatively. \citet{hunt2025} provide a weather-driven account of
Indian state demand without isolating per-festival magnitudes. \citet{banerjee2010}
is a useful foil: a Northern-grid study that reports no holiday effect, which our
weather-adjusted estimates contradict once the data are corrected
(Section~\ref{sec:data}).

\paragraph{Holiday modelling in load forecasting.} Holiday handling is a recurring
theme in the forecasting literature. \citet{ziel2018} studies the treatment of
public holidays in load models, and \citet{arora2018} incorporate special-day
effects in short-term forecasting. The general point we build on is that rare,
recurring special days are hard to encode well with dummies alone; impact-adjusted
or target-style encodings of such days are an established device, and our
signed-magnitude feature sits in that family.

\paragraph{Counterfactual demand estimation.} Estimating ``what demand would have
been'' absent a shock is methodologically close to the counterfactual estimation
used to quantify COVID-19 lockdown effects on electricity demand
\citep{fezzi2020,bahmanyar2020}. Machine-learning counterfactual and global
forecasting approaches \citep{prasanna2022,rubattu2023} are natural alternative
baselines for the robustness grid. Finally, the forecast-evaluation framing
follows standard practice on forecast comparison and encompassing
\citep{clements1996}.

\paragraph{Prior disclosure.} The signed-effect idea was previously disclosed as
part of the Gridalytics project (a public pull request and an internal research
note); we treat it here as our own prior work and extend it with the
decade-long estimation and the formal forecast ablation.

% =========================================================================
\section{Data}
\label{sec:data}

\paragraph{Demand.} The primary series is Grid-India's Power Supply Position (PSP)
daily peak demand met for Delhi, 2015-01-01 to 2024-11-20, a continuous
ten-year record. At daily resolution the target is the day's \emph{peak} (the PSP
reports peak demand met); where the series is derived from higher-resolution SCADA
data, days are aggregated with \texttt{max}, never \texttt{mean}.

\paragraph{Weather.} Daily weather is ERA5 reanalysis obtained through the
Open-Meteo archive API for Delhi ($28.61^\circ$N, $77.21^\circ$E): mean, maximum
and minimum 2\,m temperature, mean relative humidity, and precipitation sum. We
derive cooling and heating degree days,
$\mathrm{CDD}_t=\max(\bar T_t-24,0)$ and $\mathrm{HDD}_t=\max(18-\bar T_t,0)$,
their squares, and one- and two-day lags to capture the thermal inertia of the
building stock.

\paragraph{Calendar.} Festival and holiday dates are taken from the
\texttt{holidays} library for India using both the \emph{public} and
\emph{optional} categories, then mapped to a canonical, Delhi-relevant list of up
to 25 named festivals. The optional category is essential: the public category
alone omits Holi, Govardhan Puja, Bhai Duj, Chhath and Karwa Chauth.

\subsection{Data-integrity audit}
\label{sec:integrity}

An earlier pass over this pipeline concluded that festivals barely mattered and
that a magnitude feature was neutral. That conclusion was an artifact of several
silent defects, each of which nudged the result toward ``no effect.'' We document
them because the detection tests are cheap and generally reusable, and because
the corrected result (Section~\ref{sec:results1}) is the paper.

\begin{description}
\item[\textnormal{(D1)} Report-date lag (resolved).] Pre-2021 \texttt{psp\_daily}
  rows were dated one day late---a PSP report labelled day $D$ carries day $D{-}1$'s
  demand. Evidence: the weekly demand low falls on Sunday/Monday in 2015--2020 but
  Saturday/Sunday in 2021--2024, with the break near 2021-01-02 (a cumulative-sum
  test); the day-of-month pattern is identical for $\le 12$ and $>12$, so this is a
  one-day shift, not a day/month swap. Detection: a day-of-week-profile alignment
  test across sub-periods.
\item[\textnormal{(D2)} Day/month swap (resolved).] Legacy 5-minute SCADA rows
  (2021 to late 2025) had day and month transposed whenever the day was $\le 12$,
  from a date parse without \texttt{dayfirst=True}. Checking the 5-minute daily
  maximum against the PSP peak, the stored date gives median APE $24.3\pct$ and
  correlation $0.002$; the swapped date gives $0.6\pct$ and $0.999$ ($n{=}492$),
  while days $>12$ are already correct ($0.7\pct$, $0.998$). Detection: a
  cross-source agreement test at offsets $-1/0/+1$ and with day/month swapped.
\item[\textnormal{(D3)} Missing festivals (resolved).] The off-the-shelf calendar
  omitted Holi and several optional-category festivals before 2024. Detection: a
  calendar-completeness check against an optional-holiday list.
\item[\textnormal{(D4)} Peak/mean mixing (resolved).] After 2024-11 the daily
  target blended PSP peak with a 5-minute daily \emph{mean} (${\approx}0.83\times$
  peak), faking a level drop; the target is now the peak everywhere.
\item[\textnormal{(D5)} Target leak (resolved).] Differenced and rolling target
  features (\texttt{delhi\_mw\_diff\_*}, rolling statistics) included the current
  target $y_t$ during training but not at inference, flattering cross-validated
  accuracy. This is a textbook leak and we treat it as a cautionary case study in
  its own right (Section~\ref{sec:leak}).
\end{description}

Defects D1--D4 are fixed in the code, and the local database was repaired with an
idempotent script (dry run by default). D5 was subsequently fixed on the main
branch so that all differenced and rolling features end at $t{-}1$.
Table~\ref{tab:before-after} shows how the festival estimates and model fit moved
once D1--D3 were corrected.

\begin{table}[t]
\centering
\caption{Effect of the data-integrity corrections (D1--D3) on the festival
estimates and the event-study fit. The ``before'' column is the earlier pass on
broken, weather-unadjusted data; the ``after'' column is the corrected,
weather-adjusted estimate used throughout this paper.}
\label{tab:before-after}
\begin{tabular}{@{}lrr@{}}
\toprule
Quantity & Before (broken data) & After (corrected) \\
\midrule
Diwali effect        & $-10.5\pct$ (shrunk) & $-17.2\pct$ \\
Holi effect          & absent (not in calendar) & $-31.2\pct$ \\
Republic Day effect  & $-3.9\pct$  & $-14.2\pct$ \\
Id-ul-Fitr effect    & $+3.0\pct$ (sign artifact) & ${\approx}\,0$ (n.s.) \\
Magnitude feature    & neutral & helps on festival days \\
\midrule
Model $R^2$          & $0.953$ & $0.967$ \\
Residual SD          & $5.21\pct$ & $4.39\pct$ \\
\bottomrule
\end{tabular}
\end{table}

The improvement in fit ($R^2$ $0.953\to0.967$; residual standard deviation
$5.21\pct\to4.39\pct$) is independent evidence that the one-day shift was real:
correcting it makes the weather model fit better, not worse. The corrupted dates
also explain a misleading earlier ``per-year'' pattern, in which Diwali and Holi
appeared to have near-zero effects in 2015--2020 and large effects only in
2021--2024.

% =========================================================================
\section{Method}
\label{sec:method}

\subsection{Event-study counterfactual}

We model the log daily peak as a function of weather, calendar fixed effects,
lockdown controls and festival indicators:
\begin{align}
\log L_t
  &= \alpha
   + g(\text{weather}_t)
   + \sum_{d} \gamma_d\,\mathrm{DOW}_{d,t}
   + \sum_{m} \delta_m\,\mathrm{YM}_{m,t} \nonumber\\
  &\quad + \kappa_1\,\mathrm{COVID1}_t
   + \kappa_2\,\mathrm{COVID2}_t
   + \sum_{e} \beta_e\,F_{e,t}
   + \varepsilon_t ,
\label{eq:eventstudy}
\end{align}
where $\mathrm{DOW}$ are day-of-week dummies, $\mathrm{YM}$ are year-month fixed
effects (absorbing slow seasonal and trend movement), and $\mathrm{COVID1}$,
$\mathrm{COVID2}$ flag the two national lockdown windows (2020-03-23 to
2020-05-31 and 2021-04-19 to 2021-06-07). The weather term $g(\cdot)$ comprises
$\mathrm{CDD}$, $\mathrm{CDD}^2$, $\mathrm{HDD}$, $\mathrm{HDD}^2$, their one- and
two-day lags, maximum temperature, mean relative humidity, $\log(1+\text{precip})$,
and a $\mathrm{CDD}\times\mathrm{RH}$ interaction. Each $F_{e,t}$ is an indicator
for festival $e$ (those with at least three occurrences in the sample are
retained, giving up to 25 festival dummies). The reported effect of festival $e$
is the percentage deviation from the weather-adjusted counterfactual,
$\exp(\beta_e)-1$. The model is estimated by OLS on $N=\num{3569}$ days with
Newey--West HAC standard errors at lag~7, so that serial correlation and
heteroskedasticity over the weekly cycle are accommodated.

\subsection{Inference: placebo randomisation}
\label{sec:placebo}

HAC standard errors turn out to be anti-conservative here. We run a placebo test
that shifts every festival date by $\pm 28$ and $\pm 35$ days (keeping the same
day of week) and re-estimates Equation~\eqref{eq:eventstudy}. The placebo
coefficients have mean $0.5\pct$ and standard deviation $1.9\pct$, as they should
under the null; but $22\pct$ of them reach $p<0.05$ against a nominal $5\pct$,
confirming that the HAC $p$-values overstate significance. We therefore read
significance against the placebo distribution rather than the HAC $p$-values: the
top seven effects lie more than three placebo standard deviations from zero and
survive regardless. The pre-submission inference will move fully to
randomisation (placebo) or block-bootstrap intervals
(Section~\ref{sec:limits}).

\subsection{Multi-day windows and per-year stability}

To characterise structure beyond the festival day itself, we fit the
counterfactual on non-event days only, predict it for every day, and read the
deviation $(\exp(\log L_t - \widehat{\log L_t})-1)$ over a window around each
festival ($d{-}3$ to $d{+}3$ for Diwali, $d{-}2$ to $d{+}2$ for Holi). The same
per-day deviations, grouped by year, give the year-to-year stability profile.

\subsection{Forecast ablation}
\label{sec:method-ablation}

To test operational value we run a rolling-origin day-ahead backtest: for each
test year $Y \in \{2019,\dots,2024\}$ we train on all earlier years and forecast
$Y$. A LightGBM model (600 trees, learning rate $0.03$, 31 leaves) predicts the
log peak from weather, calendar fields, autoregressive lags and a seven-day
rolling mean, at horizons $h=1$ (latest usable lag $D{-}1$) and $h=7$ (latest lag
$D{-}7$). Weather is supplied as an \emph{oracle} (actual ERA5), which upper-bounds
the achievable accuracy; replacing it with archived forecasts is a pre-submission
task (Section~\ref{sec:limits}). We compare five festival-information encodings:

\begin{description}
\item[none] no holiday information;
\item[DB flag] the binary flag as shipped in the repository calendar (with Holi
  missing before 2024);
\item[full flag] a binary flag from the complete public$+$optional calendar;
\item[flag $+$ identity] the full flag plus a categorical festival identity;
\item[flag $+$ magnitude] the full flag plus a causal signed magnitude.
\end{description}

The signed magnitude for a festival on a given date is the mean deviation of that
festival's \emph{strictly prior-year} occurrences from a weather/calendar
counterfactual; a never-before-seen occurrence gets zero. The feature is fully
out-of-sample within each fold: for test year $Y$ the counterfactual is itself
fit only on non-event days \emph{strictly before} $Y$, with a continuous year
trend replacing year fixed effects so it extrapolates cleanly to the held-out
year. We report MAPE separately on the 147 festival test days and on normal days,
against a persistence baseline, and compare encodings on festival days with a
paired Wilcoxon signed-rank test.

% =========================================================================
\section{Results I: the festival catalogue}
\label{sec:results1}

\subsection{Per-festival effects}

Table~\ref{tab:effects} reports the weather-adjusted effect of each festival on
Delhi's daily peak, with $95\pct$ HAC confidence intervals.
Figure~\ref{fig:forest} shows the same effects as a diverging forest plot. The
sign varies: most festivals suppress the peak---offices and industry shut---while
a few raise it. Holi is the single largest effect at $-31.2\pct$. The Diwali
cluster (Diwali, Govardhan Puja, Bhai Duj) and Republic Day are all large and
tightly estimated. Eid-ul-Fitr/Bakrid is statistically null---notable because the
earlier broken-data pass had reported a spurious $+3\pct$. Christmas and Maha
Shivaratri are positive but flagged for confounds (a winter/New-Year overlap and a
February--March weather transition respectively) and are reported as provisional.

\begin{table}[t]
\centering
\caption{Event-study effect of each festival on Delhi's daily peak demand versus a
weather-adjusted counterfactual, 2015--2024. Intervals are $95\pct$ HAC(7).
``Years'' is the number of in-sample occurrences. Effects are ordered from largest
suppression to largest increase. The top seven survive placebo-randomisation
inference by more than three placebo standard deviations
(Section~\ref{sec:placebo}). \emph{Draft estimates}; to be regenerated from
\texttt{scripts/paper/festival\_eventstudy.py} with the final weather series.}
\label{tab:effects}
\begin{tabular}{@{}l r l c@{}}
\toprule
Festival & Effect & $95\pct$ CI & Years \\
\midrule
Holi                       & $-31.2\pct$ & $[-32.6,\,-29.8]$ & 10 \\
Govardhan Puja (Diwali$+1$) & $-18.7\pct$ & $[-22.1,\,-15.3]$ & 10 \\
Diwali                     & $-17.2\pct$ & $[-19.4,\,-14.9]$ & 10 \\
Republic Day               & $-14.2\pct$ & $[-16.3,\,-12.1]$ & 10 \\
Bhai Duj (Diwali$+2$)      & $-11.7\pct$ & $[-14.5,\,-8.9]$  & 10 \\
Dussehra                   & $-10.7\pct$ & $[-14.0,\,-7.3]$  & 10 \\
Ram Navami                 & $-6.2\pct$  & $[-8.1,\,-4.3]$   & 10 \\
Independence Day           & $-3.8\pct$  & $[-6.5,\,-1.1]$   & 8 \\
Eid-ul-Fitr / Bakrid       & ${\approx}\,0$ & n.s. & 10 \\
Christmas$^\dagger$        & $+4.1\pct$  & $[+2.4,\,+5.9]$   & 9 \\
Maha Shivaratri$^\dagger$  & $+6.0\pct$  & $[+3.5,\,+8.4]$   & 9 \\
\bottomrule
\end{tabular}

\vspace{2pt}
{\footnotesize $^\dagger$ Positive sign provisional: Christmas carries a
winter/New-Year confound; Maha Shivaratri falls in the February--March weather
transition. To be explained or reported as unexplained before submission.}
\end{table}

% TODO: regenerate figures/festival_effects_forest.pdf from scripts/paper/festival_eventstudy.py
\begin{figure}[t]
\centering
\placeholderfig{festival_effects_forest}{5.5cm}
\caption{Per-festival effect on Delhi's daily peak demand with $95\pct$
intervals (diverging forest plot of Table~\ref{tab:effects}).}
\label{fig:forest}
\end{figure}

\subsection{Per-year stability}

The large effects are stable from year to year, not driven by any single
occurrence. Holi falls between $-27\pct$ and $-34\pct$ in \emph{every} year of
2015--2024. Diwali falls between $-11\pct$ and $-24\pct$, and the drop is
deepening---the four largest Diwali dips are 2021--2024.
Figure~\ref{fig:stability} plots the per-year deviations for Holi and Diwali.

% TODO: regenerate figures/peryear_stability.pdf from scripts/paper/festival_eventstudy.py
\begin{figure}[t]
\centering
\placeholderfig{peryear_stability}{5cm}
\caption{Per-year deviation of Holi and Diwali from the weather-adjusted
counterfactual, 2015--2024. Holi is stable within $[-34,-27]\pct$; the Diwali dip
deepens over the decade.}
\label{fig:stability}
\end{figure}

\subsection{Multi-day structure}

Festivals are not one-day events. The Diwali effect is a three-day trough:
Diwali ${\approx}-17\pct$, Govardhan Puja (the following day) ${\approx}-19\pct$,
and Bhai Duj (two days after) ${\approx}-12\pct$. Holi carries a day-after effect
of about $-13\pct$ in addition to the $-31\pct$ on Holi itself.
Figure~\ref{fig:windows} shows both windows. A single-day binary flag cannot
express either shape, and even a per-festival identity misses the ordered,
multi-day decay---motivating the signed, continuous encoding tested next.

% TODO: regenerate figures/event_windows.pdf from scripts/paper/festival_eventstudy.py
\begin{figure}[t]
\centering
\placeholderfig{event_windows}{5cm}
\caption{Event windows: mean deviation from the counterfactual around Diwali
($d{-}3$ to $d{+}3$) and Holi ($d{-}2$ to $d{+}2$). Diwali is a three-day trough;
Holi carries a day-after effect.}
\label{fig:windows}
\end{figure}

% =========================================================================
\section{Results II: forecasting value}
\label{sec:results2}

Table~\ref{tab:ablation} reports the rolling-origin backtest. On festival days the
signed magnitude is the clear winner at both horizons; the ordering of the other
encodings is informative. Moving from no information to a binary flag helps a
little; using the \emph{complete} calendar flag (adding Holi and the optional
festivals) helps more; but adding a categorical festival identity on top of the
flag does \emph{not} improve on the flag, and the signed magnitude is what closes
the gap.

\begin{table}[t]
\centering
\caption{Day-ahead festival-information ablation. MAPE ($\pct$) on test years
2019--2024 (rolling origin), LightGBM forecasting the log peak, oracle weather,
147 festival test days. The flag$+$magnitude counterfactual is \emph{fully
out-of-sample} (for each test year it is fit only on non-event days strictly
before that year). Best festival-day entry in each horizon in bold. Weather is
still oracle; archived-forecast weather, Diebold--Mariano tests and multiple
seeds remain pre-submission tasks (Section~\ref{sec:limits}).}
\label{tab:ablation}
\begin{tabular}{@{}l l r r r r r r@{}}
\toprule
$h$ & day type & none & DB flag & full flag & flag$+$id & flag$+$mag & persist. \\
\midrule
1 & festival & $7.94$ & $7.00$ & $6.61$ & $6.70$ & $\mathbf{5.50}$ & $7.37$ \\
1 & normal   & $4.23$ & $4.03$ & $4.06$ & $4.05$ & $4.06$          & $4.55$ \\
7 & festival & $10.33$& $9.38$ & $9.45$ & $9.50$ & $\mathbf{6.90}$ & $13.35$ \\
7 & normal   & $6.73$ & $6.49$ & $6.41$ & $6.48$ & $6.42$          & $9.21$ \\
\bottomrule
\end{tabular}
\end{table}

Concretely, on festival days the signed magnitude cuts error from the full-flag
baseline of $6.6\pct$ to $5.5\pct$ at $h=1$ (Wilcoxon $p=0.005$) and from
$9.5\pct$ to $6.9\pct$ at $h=7$ ($p\approx1\times10^{-5}$). The gain is sharply
concentrated on the festival days themselves: on normal days, where the magnitude
feature is exactly zero, every encoding moves by less than $0.3$ points, which is
how we know the improvement is the event signal and not a global perturbation of
the model. The effect is largest where the festival is largest---Holi festival-day
MAPE at $h=1$ falls to about $11\pct$, from $25\pct$ under the complete flag (and
$43\pct$ with no festival information at all).

The mechanism is a rare-event encoding one. A festival occurs about once per year,
so over a five-to-nine-year training span a gradient-boosted tree sees only a
handful of positive examples; it has nowhere near enough splits to learn a
rare event's size from a binary dummy or a bare identity. A signed, pre-computed
magnitude instead hands the model the answer as a continuous prior. This is a
feature-encoding result, in the spirit of impact-adjusted or target-style
encodings of special days \citep{ziel2018,arora2018}; it is not a new estimator.

The magnitude feature does not win uniformly. Bhai Duj at $h=1$ gets
\emph{worse} under the magnitude encoding ($3.2\pct \to 7.3\pct$). The cause is
structural: Bhai Duj falls two days after Diwali, so its $h=1$ lag is Govardhan
Puja (Diwali$+1$), the deepest trough in the catalogue ($-18.7\pct$), in all ten
years. The model then sees a deep-trough lag \emph{and} a signed-magnitude prior
pointing further down, and over-predicts the fall --- whereas demand in fact
recovers from Govardhan Puja ($-18.7\pct$) to the shallower Bhai Duj
($-11.7\pct$). Exposing the lag day's own magnitude as an extra feature, so the
model can mean-revert, helps only marginally (Bhai Duj $7.3\pct \to 7.0\pct$,
overall festival-day MAPE $5.50\pct \to 5.42\pct$) while slightly hurting
Govardhan Puja, so we do not adopt it; the proper remedy is a window- or
sequence-aware encoding of consecutive event days, which we leave to future work.

\subsection{A target leak, and why trees exploit it too}
\label{sec:leak}

Defect D5 deserves a result of its own. The differenced and rolling target
features included the current target $y_t$ in training but not at inference. On
the daily series this made the as-shipped cross-validated MAPE look like
$1.08\pct$ when the leak-free value is $3.92\pct$, against $5.05\pct$ for
persistence. After shifting all such features to end at $t{-}1$, the leak-free
champions are $3.92\pct$ (daily), $1.31\pct$ (hourly) and $0.45\pct$ (5-minute)
MAPE, versus persistence baselines of $5.05\pct$, $3.43\pct$ and $0.68\pct$. The
point for practitioners is that this is \emph{not} a linear-model-only
pathology: gradient-boosted trees exploit a target leak just as eagerly as a
linear model does, and the only reliable guard is to construct differenced and
rolling features so they end at $t{-}1$ and to verify train/serve parity. We
include it as a cautionary case study alongside the festival result, since both
are instances of the same discipline---honest numbers only after the pipeline is
audited.

% =========================================================================
\section{Discussion}
\label{sec:discussion}

\paragraph{Why binary flags fail for rare events.} The central practical lesson is
that a binary \texttt{is\_festival} flag is the wrong representation for a shock
that is large, signed, multi-day and rare. It cannot carry the sign (Christmas up,
Holi down), the magnitude (Holi is four times Dussehra), or the multi-day shape
(the Diwali trough). A categorical identity carries sign and shape in principle but
cannot be learned from one example per year. A signed, causal magnitude carries
all three and is learnable because it is pre-computed from prior years, not
rediscovered from dummies.

\paragraph{What the earlier corrupted pipeline concluded.} Before the audit, the
same data said festivals barely mattered and that the magnitude feature was
neutral (Table~\ref{tab:before-after}). Each defect pushed in the same direction,
and together they produced a publishable-looking negative result. The corrected
analysis reverses it. We think the general lesson---that a cluster of small,
silent data defects can manufacture a confident null---is worth more than the
specific Delhi numbers.

\paragraph{Operational use.} For an SLDC scheduling a day ahead, a reliable
$-17\pct$ to $-31\pct$ festival-day correction, with a three-day Diwali shape, is
directly actionable for unit commitment and reserve sizing; for a DISCOM, it sizes
day-ahead power purchase on exactly the days when a mis-sized forecast is most
costly. We are careful, however, not to oversell the forecasting lever: in the
production fleet, where rich autoregressive lags already encode most of the
one-step signal, the signed magnitude earns its place in the champion model only
for some series (clearly for Punjab's daily peak, marginally or not at all
elsewhere), and it is shipped only where a per-state back-test says it pays. The
\emph{quantification} is broad and robust; the \emph{forecasting} lever is real
but narrow.

% =========================================================================
\section{Limitations and future work}
\label{sec:limits}

\begin{description}
\item[Daily peak, one city.] The decade-long estimation is Delhi's daily peak
  only. We make no sub-daily or multi-city accuracy claim. Festivals also reshape
  the intraday profile (the Diwali evening peak), which no daily scalar can
  express; quantifying that needs a corrected 5-minute backfill around event dates
  and is the one lever with real upside left.
\item[Inference.] HAC(7) standard errors are anti-conservative here
  (Section~\ref{sec:placebo}); the final inference will use
  placebo-randomisation or block-bootstrap intervals throughout. The top seven
  effects survive this already.
\item[Oracle weather.] The forecast ablation uses actual (oracle) weather, which
  upper-bounds accuracy. Before submission we will substitute archived weather
  \emph{forecasts} (Open-Meteo historical-forecast API, 2022 onward) or a
  noise-injection sensitivity check, and add Diebold--Mariano tests and multiple
  seeds.
\item[Robustness grid.] We will report alternative counterfactuals (a $\pm21$-day
  same-day-of-week estimator, this regression, and a machine-learning
  counterfactual \citep{prasanna2022,rubattu2023}); results with the 2020--2021
  COVID years dropped; peak versus energy (re-scraping PSP energy for legacy rows);
  and the full $d{-}1$ to $d{+}3$ window.
\item[Positive effects.] Christmas and Maha Shivaratri are positive but confounded
  (winter/New-Year overlap; February--March weather transition). We will either
  explain them or report them explicitly as unexplained.
\item[Scoop risk.] We have not yet checked the state-level Grid-India demand
  pattern tables (Vol.~II) for any Delhi festival magnitudes; this is the main
  residual novelty risk and will be resolved before we claim ``first.''
\end{description}

% =========================================================================
\section*{Reproducibility}

The event-study estimator and the forecast ablation are implemented in
\texttt{scripts/paper/festival\_eventstudy.py} and
\texttt{scripts/paper/festival\_forecast\_ablation.py}. The data-repair script and
the leak-free feature construction are part of the Gridalytics repository.
Weather is fetched from the Open-Meteo archive API. The tables and figures in this
draft are the feasibility estimates of 2026-10-29 and are to be regenerated with
the final weather series before submission.

% =========================================================================
\bibliographystyle{plainnat}
\bibliography{references}

\end{document}
